\section{Related work}
\label{sec:related}

\subsection{Desktop grids and volunteer computing}

Harvesting idle cycles from machines whose owners have first claim is a mature
field. BOINC established the platform model, in which participants run a client
that pulls work units and the scheduling model belongs to the platform; later
systems such as Pando pursue the same goal in browsers with a declarative
concurrent programming model. Intermittency and owner control are founding
assumptions throughout, and the field has developed a substantial literature on
tolerating them.

The difference in emphasis is where the intermittency is handled. That
literature addresses it at the \emph{application} layer, through checkpointing,
replication and work-unit redundancy, and treats the network as a solved
substrate because participants reach a central server outbound. This note
addresses a case that model does not encounter: execute nodes that must be
reachable \emph{inbound} while moving between subnets that cannot route to each
other (\S\ref{sec:mobility}).

\subsection{Requested versus actual resource usage}

That jobs request far more than they use, and that this depresses utilisation,
is well established in datacentre practice; published cluster traces and the
right-sizing systems built on them treat over-request as the normal case rather
than an anomaly. The usual remedy is to correct the \emph{request}, by
recommending or automatically adjusting it from observed history.

The finding in \S\ref{sec:ledger} is adjacent but distinct. The request is left
alone, and the \emph{ledger derived from it} is replaced with a measurement, on
the grounds that the reservation the ledger enforces is not being honoured in
the first place. That is a narrower claim, and it is specific to a system whose
per-job enforcement has been deliberately disabled, but it removes a gate that
sits upstream of policy rather than adjusting an input to it.

\subsection{Pressure stall information}

PSI is a Linux kernel facility for quantifying time lost to resource contention,
and its use in scheduling is established outside batch systems: container
orchestrators expose it for node-level decisions, and Android's low-memory
daemon uses it in place of an in-kernel killer precisely because it reacts
faster than free-memory thresholds on constrained devices. The observation that
a scheduler should prefer nodes with lower \emph{pressure} rather than lower
\emph{utilisation} is not new.

What appears not to be written down is its use inside a batch system's own
policy expressions, where it competes with the load-average signal those systems
have always used, and where the ten-second reaction window matters because a
negotiation cycle can admit an entire burst inside one decay constant of a
one-minute mean (\S\ref{sec:admission}).

\subsection{HTCondor networking}

The multi-homing mechanism this note relies on is documented: a daemon may
advertise a public and a private address, and peers select between them by
matching a configured private network name. The documentation is explicit that
two addresses is the limit.

What we did not find described is the consequence for machines that \emph{move}.
Because the selection is made by the reader, a configured value is silently
wrong the moment a machine changes segment, and the failure is not a refusal but
a job that runs to completion and then cannot return its output
(\S\ref{sec:mobility}). Deriving the value at startup from the collector's own
advertisement, so that no machine holds a network fact, is the part of this work
we believe to be unreported.
